% EGE 2027, task 18. Source IDs and URLs are kept in comments.
% Editable task bank. Append new BankTask blocks at the end of the document.
\documentclass[12pt,a4paper]{article}
\usepackage{physics-textbook}
\usepackage{graphicx,needspace}
\newif\ifShowSolutions
\ifdefined\TutorEdition
  \ShowSolutionstrue
\else
  \ShowSolutionsfalse
\fi
\newcounter{banktask}
\newcommand{\BankTask}[3]{%
  \refstepcounter{banktask}%
  \Needspace{35mm}\par\medskip\noindent\begin{minipage}{\linewidth}%
  \noindent\textbf{\thebanktask.}\quad #2\par
  \ifShowSolutions
    #3\par
    \ifdefined\TutorEdition
      \ifcsname TutorSource#1\endcsname
        \smallskip{\small\csname TutorSource#1\endcsname\par}%
      \fi
    \fi
  \fi
  \end{minipage}\par\medskip
}
\begin{document}


% Source 0FDA4F: https://ege.fipi.ru/bank/questions.php?search=1&proj=BA1F39653304A5B041B656915DC36B38&qid=0FDA4F
\BankTask{0FDA4F}{
\begin{enumerate}
\item Тела действуют друг на друга силами одной и той же природы, равными по модулю и противоположными по направлению.
\item Конденсацией называют процесс преобразования жидкости в твёрдое вещество.
\item В процессе поляризации связанные положительные и отрицательные заряды диэлектрика смещаются в противоположные стороны.
\item Энергия магнитного поля катушки индуктивностью \(L\) прямо пропорциональна квадрату силы тока в катушке.
\item Свет обладает дуализмом свойств: при его распространении проявляются корпускулярные свойства света, а при взаимодействии с веществом -- волновые.
\end{enumerate}
}{
\smallskip\noindent\textbf{Решение от GPT}\par
1) Верно по третьему закону Ньютона: \(\vec F_{12}=-\vec F_{21}\), силы имеют одну природу. 2) Неверно: конденсация — переход из газообразного состояния в жидкое. 3) Верно: при поляризации связанные заряды противоположных знаков смещаются в противоположные стороны. 4) Верно, поскольку \(W=LI^2/2\). 5) Неверно: при распространении света проявляются волновые свойства, а при поглощении и испускании — корпускулярные. Ответ: \(1,\,3,\,4\).
}

% Source A1CBB8: https://ege.fipi.ru/bank/questions.php?search=1&proj=BA1F39653304A5B041B656915DC36B38&qid=A1CBB8
\BankTask{A1CBB8}{
\begin{enumerate}
\item Модуль сил гравитационного взаимодействия двух материальных точек обратно пропорционален квадрату расстояния между ними.
\item Давление насыщенного пара увеличивается с ростом абсолютной температуры пара и не зависит от его объёма.
\item В однородном электростатическом поле работа силы электростатического поля по перемещению заряда между двумя точками прямо пропорциональна длине траектории.
\item При переходе электромагнитной волны из оптически менее плотной в оптически более плотную среду частота волны остаётся неизменной.
\item При распространении света проявляются только его корпускулярные свойства, а при взаимодействии с веществом -- только волновые.
\end{enumerate}
}{
\smallskip\noindent\textbf{Решение от GPT}\par
1) Верно: \(F=Gm_1m_2/r^2\). 2) Верно: давление насыщенного пара растёт с температурой и при данной температуре не зависит от объёма, пока пар остаётся насыщенным. 3) Неверно: \(A=q(\varphi_1-\varphi_2)\), работа зависит от конечных точек, а не от длины траектории. 4) Верно: при переходе через границу сред частота сохраняется, \(\nu_1=\nu_2\). 5) Неверно: свет проявляет волновые свойства при интерференции и дифракции, а корпускулярные — например, при фотоэффекте. Ответ: \(1,\,2,\,4\).
}

% Source 024D10: https://ege.fipi.ru/bank/questions.php?search=1&proj=BA1F39653304A5B041B656915DC36B38&qid=024D10
\BankTask{024D10}{
\begin{enumerate}
\item Громкость звука определяется частотой колебаний.
\item Температура плавления кристаллических тел зависит от их массы.
\item В цепи постоянного тока на всех параллельно соединённых резисторах напряжение одинаково.
\item Скорость распространения радиоволн в вакууме равна скорости света в вакууме.
\item При \(\beta\)-распаде ядра образуется ион нового элемента и ядро атома гелия.
\end{enumerate}
}{
\smallskip\noindent\textbf{Решение от GPT}\par
1) Неверно: частота определяет высоту звука, а громкость связана прежде всего с интенсивностью звука. 2) Неверно: при одинаковом давлении температура плавления данного кристаллического вещества не зависит от массы. 3) Верно: параллельно соединённые резисторы подключены к одним и тем же двум узлам, поэтому \(U_1=U_2=\dots\). 4) Верно: радиоволны являются электромагнитными волнами и распространяются в вакууме со скоростью \(c=3\cdot10^8\,\text{м/с}\). 5) Неверно: при \(\beta\)-распаде испускается электрон или позитрон и соответствующее антинейтрино или нейтрино, а ядро гелия испускается при \(\alpha\)-распаде. Ответ: \(3,\,4\).
}

% Source 46BA23: https://ege.fipi.ru/bank/questions.php?search=1&proj=BA1F39653304A5B041B656915DC36B38&qid=46BA23
\BankTask{46BA23}{
\begin{enumerate}
\item[1)] Давление столба жидкости на дно сосуда прямо пропорционально её плотности.
\item[2)] Удельная теплота плавления вещества показывает, какое количество теплоты необходимо сообщить \(1\,\text{кг}\) вещества, находящемуся при температуре плавления, чтобы его расплавить.
\item[3)] В процессе электризации трением два первоначально незаряженных тела приобретают разноимённые и различные по модулю заряды.
\item[4)] При переходе света из оптически более плотной среды в оптически менее плотную среду угол падения больше угла преломления.
\item[5)] При \(\alpha\)-распаде ядра выполняются закон сохранения электрического заряда, закон сохранения импульса.
\end{enumerate}
}{
\smallskip\noindent\textbf{Решение от GPT}\par
1) Верно: гидростатическое давление \(p=\rho gh\) прямо пропорционально плотности жидкости. 2) Верно: из \(Q=\lambda m\) следует, что удельная теплота плавления равна количеству теплоты для плавления \(1\,\text{кг}\) вещества при температуре плавления. 3) Неверно: по закону сохранения заряда \(q_1+q_2=0\), поэтому заряды равны по модулю. 4) Неверно: по закону преломления \(n_1\sin\alpha=n_2\sin\beta\) при \(n_1>n_2\) угол преломления больше угла падения при наклонном падении и наличии преломлённого луча. 5) Верно: при \(\alpha\)-распаде сохраняются электрический заряд и импульс; ответ: \(1,\ 2,\ 5\).
}

% Source 027C91: https://ege.fipi.ru/bank/questions.php?search=1&proj=BA1F39653304A5B041B656915DC36B38&qid=027C91
\BankTask{027C91}{
\begin{enumerate}
\item[1)] При увеличении длины нити математического маятника период его колебаний уменьшается.
\item[2)] Явление диффузии протекает в твёрдых телах значительно медленнее, чем в жидкостях.
\item[3)] Сила Лоренца отклоняет положительно и отрицательно заряженные частицы, влетающие под углом к линиям индукции однородного магнитного поля, в противоположные стороны.
\item[4)] Дифракция рентгеновских лучей невозможна.
\item[5)] В процессе фотоэффекта с поверхности вещества под действием падающего света вылетают электроны.
\end{enumerate}
}{
\smallskip\noindent\textbf{Решение от GPT}\par
1) Неверно: \(T=2\pi\sqrt{\ell/g}\), поэтому с увеличением длины период растёт. 2) Верно: диффузия в твёрдых телах значительно медленнее вследствие меньшей подвижности частиц. 3) Верно: сила \(\vec F=q\,\vec v\times\vec B\) при смене знака заряда меняет направление на противоположное. 4) Неверно: рентгеновские лучи дифрагируют, например, на кристаллической решётке. 5) Верно: внешний фотоэффект представляет собой вылет электронов из вещества под действием света; ответ: \(2,\ 3,\ 5\).
}

% Source 5A149B: https://ege.fipi.ru/bank/questions.php?search=1&proj=BA1F39653304A5B041B656915DC36B38&qid=5A149B
\BankTask{5A149B}{
\begin{enumerate}
\item[1)] Сила трения скольжения – сила гравитационной природы.
\item[2)] Для конденсации жидкости ей необходимо сообщить положительное количество теплоты.
\item[3)] Силой Лоренца называют силу, с которой однородное электрическое поле действует на постоянные магниты.
\item[4)] Линейчатый спектр дают вещества в газообразном атомарном состоянии.
\item[5)] Количество фотоэлектронов, вылетающих с поверхности металла за единицу времени, прямо пропорционально интенсивности падающего на поверхность металла света.
\end{enumerate}
}{
\smallskip\noindent\textbf{Решение от GPT}\par
1) Неверно: сила трения имеет электромагнитную природу. 2) Неверно: при конденсации пар отдаёт теплоту, \(Q=-Lm\). 3) Неверно: магнитная сила Лоренца действует на движущийся заряд, \(\vec F=q\,\vec v\times\vec B\). 4) Верно: переходы между дискретными энергетическими уровнями атомов дают линейчатый спектр. 5) Верно: при неизменной частоте света выше красной границы число вылетающих фотоэлектронов за единицу времени прямо пропорционально интенсивности света; ответ: \(4,\ 5\).
}

% Source 777063: https://ege.fipi.ru/bank/questions.php?search=1&proj=BA1F39653304A5B041B656915DC36B38&qid=777063
\BankTask{777063}{
\begin{enumerate}
\item При прохождении математическим маятником положения равновесия центростремительное ускорение его груза максимально.
\item Удельная теплоёмкость вещества показывает, какое количество теплоты необходимо сообщить \(1\,\text{кг}\) вещества для его плавления.
\item При помещении проводника в электростатическое поле наблюдается явление электростатической индукции.
\item При преломлении света, падающего из среды с меньшим показателем преломления в среду с бóльшим показателем преломления, угол падения меньше угла преломления.
\item При \(\beta\)-распаде ядра выполняются законы сохранения энергии и электрического заряда, но не выполняется закон сохранения импульса.
\end{enumerate}
}{
\smallskip\noindent\textbf{Решение от GPT}\par
1) Верно: в положении равновесия скорость максимальна, поэтому максимально и \(a_{\text{ц}}=v^2/\ell\). 2) Неверно: удельная теплоёмкость определяется соотношением \(c=Q/(m\Delta T)\), а для плавления используют удельную теплоту плавления \(\lambda=Q/m\). 3) Верно: внешнее электростатическое поле вызывает перераспределение свободных зарядов проводника. 4) Неверно: по закону \(n_1\sin\alpha=n_2\sin\beta\) при \(n_2>n_1\) получаем \(\beta<\alpha\) для ненулевого угла падения. 5) Неверно: при \(\beta\)-распаде сохраняется также полный импульс всех участников процесса, включая нейтрино или антинейтрино; ответ: \(1,\ 3\).
}

% Source 4AE181: https://ege.fipi.ru/bank/questions.php?search=1&proj=BA1F39653304A5B041B656915DC36B38&qid=4AE181
\BankTask{4AE181}{
\begin{enumerate}
\item При резонансе в механической колебательной системе амплитуда установившихся вынужденных колебаний резко уменьшается.
\item Конденсацией называют процесс преобразования пара в твёрдое вещество, минуя жидкую фазу.
\item При электрическом разряде в газе перенос заряда обеспечивается только положительно заряженными ионами.
\item Вынужденными электромагнитными колебаниями называют колебания в цепи под действием внешней периодически изменяющейся электродвижущей силы.
\item В ядерных реакторах для получения энергии используются экзотермические реакции распада тяжёлых ядер.
\end{enumerate}
}{
\smallskip\noindent\textbf{Решение от GPT}\par
1) Неверно: при резонансе амплитуда установившихся вынужденных колебаний резко возрастает. 2) Неверно: конденсация означает переход пара в жидкость, а непосредственный переход в твёрдое состояние называется десублимацией. 3) Неверно: заряд в газе переносят также электроны и отрицательные ионы. 4) Верно: внешняя периодическая ЭДС поддерживает вынужденные электромагнитные колебания. 5) Верно: при делении тяжёлых ядер в реакторе выделяется энергия \(Q=\Delta m c^2>0\); ответ: \(4,\ 5\).
}

% Source 51433E: https://ege.fipi.ru/bank/questions.php?search=1&proj=BA1F39653304A5B041B656915DC36B38&qid=51433E
\BankTask{51433E}{
\begin{enumerate}
\item При равноускоренном движении ускорение тела за любые равные промежутки времени изменяется одинаково.
\item В процессе кипения жидкости при постоянном внешнем давлении её температура не меняется.
\item Сила тока короткого замыкания определяется только внутренним сопротивлением источника.
\item В поперечной механической волне колебания частиц происходят в направлении, перпендикулярном направлению распространения волны.
\item В результате \(\alpha\)-распада элемент смещается в Периодической системе элементов Д.И. Менделеева на две клетки ближе к концу.
\end{enumerate}
}{
\smallskip\noindent\textbf{Решение от GPT}\par
1) Неверно: при равноускоренном движении ускорение постоянно, \(\vec a=\mathrm{const}\), а за равные промежутки времени одинаково изменяется скорость. 2) Верно: при постоянном внешнем давлении температура кипящей жидкости остаётся постоянной. 3) Неверно: \(I_{\text{кз}}=\mathcal E/r\), поэтому ток зависит и от ЭДС источника, и от его внутреннего сопротивления. 4) Верно: в поперечной волне частицы колеблются перпендикулярно направлению её распространения. 5) Неверно: при \(\alpha\)-распаде \(Z\to Z-2\), поэтому элемент смещается на две клетки к началу таблицы; ответ: \(2,\ 4\).
}

% Source 368BED: https://ege.fipi.ru/bank/questions.php?search=1&proj=BA1F39653304A5B041B656915DC36B38&qid=368BED
\BankTask{368BED}{
\begin{enumerate}
\item[1)] Импульсом силы называется величина, равная произведению массы тела на его ускорение.
\item[2)] В изотермическом процессе для постоянной массы газа отношение объёма газа к его давлению остаётся постоянным.
\item[3)] Модуль сил взаимодействия двух точечных неподвижных заряженных тел обратно пропорционален квадрату расстояния между заряженными телами.
\item[4)] Период свободных электромагнитных колебаний в идеальном колебательном контуре увеличивается прямо пропорционально увеличению электроёмкости конденсатора.
\item[5)] В планетарной модели атома число протонов в ядре равно числу электронов в электронной оболочке нейтрального атома.
\end{enumerate}
}{
\smallskip\noindent\textbf{Решение от GPT}\par
1) Неверно: импульс постоянной силы равен $\vec F\Delta t$, а $m\vec a=\vec F$. 2) Неверно: при изотермическом процессе $pV=\mathrm{const}$. 3) Верно: по закону Кулона $F=k|q_1q_2|/r^2$. 4) Неверно: $T=2\pi\sqrt{LC}$, поэтому при постоянной индуктивности $T\propto\sqrt C$. 5) Верно: нейтральность атома означает равенство числа протонов и электронов; ответ: $35$.
}

% Source AD5973: https://ege.fipi.ru/bank/questions.php?search=1&proj=BA1F39653304A5B041B656915DC36B38&qid=AD5973
\BankTask{AD5973}{
Энергия характеризует способность тела совершать работу.

В цилиндре под поршнем расширение газа в ходе адиабатного процесса сопровождается понижением его температуры.

Если электрический ток протекает по алюминиевому проводнику, то ни при каких условиях не может наблюдаться действие тока на магнитную стрелку.

При изменении магнитного потока через площадку, охваченную замкнутым проводящим контуром, магнитное поле индукционного тока в контуре всегда увеличивает магнитный поток через эту площадку.

При $\beta$-распаде ядра выполняется закон сохранения электрического заряда.
}{
\smallskip\noindent\textbf{Решение от GPT}\par
1) Верно: энергия характеризует способность тела совершать работу. 2) Верно: при адиабатном расширении $Q=0$, газ совершает положительную работу, поэтому $\Delta U=-A<0$ и температура понижается. 3) Неверно: ток в алюминиевом проводнике создаёт магнитное поле, которое может действовать на магнитную стрелку. 4) Неверно: по правилу Ленца индукционный ток противодействует изменению магнитного потока, поэтому его поле может как увеличивать, так и уменьшать поток. 5) Верно: при $\beta$-распаде суммарный электрический заряд сохраняется; ответ: $125$.
}

% Source C2CF4D: https://ege.fipi.ru/bank/questions.php?search=1&proj=BA1F39653304A5B041B656915DC36B38&qid=C2CF4D
\BankTask{C2CF4D}{
\begin{enumerate}
\item[1)] При равномерном прямолинейном движении за любые равные промежутки времени тело совершает одинаковые перемещения.
\item[2)] В процессе изохорного нагревания постоянной массы газа он совершает положительную работу.
\item[3)] В цепи постоянного тока при последовательном соединении резисторов чем больше сопротивление резистора, тем меньше напряжение на нём.
\item[4)] Скорость распространения инфракрасного излучения в вакууме равна скорости света в вакууме.
\item[5)] Периодом полураспада называют промежуток времени, в течение которого распадается половина исходно большого числа атомов данного радиоактивного элемента.
\end{enumerate}
}{
\smallskip\noindent\textbf{Решение от GPT}\par
1) Верно: при постоянной скорости $\Delta\vec r=\vec v\Delta t$, поэтому равным промежуткам времени соответствуют одинаковые перемещения. 2) Неверно: при постоянном объёме работа газа $A=0$. 3) Неверно: при последовательном соединении ток одинаков, поэтому по закону Ома $U=IR$ напряжение больше на резисторе с большим сопротивлением. 4) Верно: инфракрасное излучение является электромагнитным и распространяется в вакууме со скоростью $c$. 5) Верно: за период полураспада число нераспавшихся атомов уменьшается вдвое; ответ: $145$.
}

% Source eeAe26: https://ege.fipi.ru/bank/questions.php?search=1&proj=BA1F39653304A5B041B656915DC36B38&qid=eeAe26
\BankTask{eeAe26}{
1) Если скорость тела с течением времени увеличивается, то вектор ускорения направлен противоположно вектору скорости тела.
\par
2) В ходе процесса кристаллизации жидкости внутренняя энергия системы «жидкость + твёрдое тело» уменьшается, а температура этой системы остаётся неизменной.
\par
3) Энергия электрического поля конденсатора, заряженного и отключённого от источника, не меняется при изменении электроёмкости конденсатора.
\par
4) При переходе электромагнитной волны из оптически менее плотной в оптически более плотную среду длина волны уменьшается.
\par
5) В планетарной модели атома число нейтронов в ядре равно числу электронов в электронной оболочке нейтрального атома.
}{
\smallskip\noindent\textbf{Решение от GPT}\par
1) Неверно: при увеличении модуля скорости \(\vec a\cdot\vec v>0\), поэтому ускорение не может быть направлено противоположно скорости.
2) Верно: при кристаллизации выделяется теплота, внутренняя энергия уменьшается, а температура остаётся постоянной.
3) Неверно: у отключённого конденсатора заряд \(q\) сохраняется, поэтому энергия \(W=q^2/(2C)\) изменяется при изменении \(C\).
4) Верно: частота сохраняется, скорость распространения уменьшается, поэтому длина волны \(\lambda=v/\nu\) уменьшается.
5) Неверно: в нейтральном атоме число электронов равно числу протонов, а не нейтронов; ответ: \(2,\ 4\).
}

% Source D92DAA: https://ege.fipi.ru/bank/questions.php?search=1&proj=BA1F39653304A5B041B656915DC36B38&qid=D92DAA
\BankTask{D92DAA}{
1) Если сила направлена перпендикулярно направлению движения тела, то работа этой силы равна нулю.
\par
2) В ходе процесса кристаллизации твёрдого тела его температура не меняется, а внутренняя энергия системы «жидкость + твёрдое тело» возрастает.
\par
3) Силы взаимодействия двух неподвижных точечных заряженных тел равны по модулю и направлены вдоль прямой, соединяющей эти тела.
\par
4) Энергия магнитного поля катушки индуктивностью \(L\) прямо пропорциональна квадрату силы тока в катушке.
\par
5) Атом излучает свет при переходе из стационарного состояния с меньшей энергией в стационарное состояние с большей энергией.
}{
\smallskip\noindent\textbf{Решение от GPT}\par
1) Верно: при перпендикулярности силы и перемещения \(A=Fs\cos 90^\circ=0\).
2) Неверно: при кристаллизации жидкости внутренняя энергия системы «жидкость + твёрдое тело» уменьшается.
3) Верно: кулоновские силы направлены вдоль соединяющей заряды прямой и по третьему закону Ньютона равны по модулю.
4) Верно: энергия магнитного поля катушки \(W=LI^2/2\).
5) Неверно: излучение происходит при переходе на уровень с меньшей энергией, причём \(h\nu=E_{\text{нач}}-E_{\text{кон}}>0\); ответ: \(1,\ 3,\ 4\).
}

% Source 970E28: https://ege.fipi.ru/bank/questions.php?search=1&proj=BA1F39653304A5B041B656915DC36B38&qid=970E28
\BankTask{970E28}{
1) При совпадении частоты вынуждающей силы и собственной частоты колебательной системы наблюдается явление резонанса.
\par
2) Процесс передачи количества теплоты от более нагретого тела к менее нагретому является обратимым.
\par
3) В замкнутой системе алгебраическая сумма зарядов всех частиц всегда равна нулю.
\par
4) Дифракция волн хорошо наблюдается в тех случаях, когда размеры препятствий меньше длины волны или сравнимы с ней.
\par
5) В планетарной модели атома в центре атома находится положительно заряженное ядро.
}{
\smallskip\noindent\textbf{Решение от GPT}\par
1) Верно: при совпадении частоты вынуждающей силы с собственной частотой системы возникает резонанс.
2) Неверно: самопроизвольная передача теплоты от более нагретого тела к менее нагретому необратима.
3) Неверно: в замкнутой системе суммарный заряд сохраняется, \(\sum_i q_i=\mathrm{const}\), но не обязательно равен нулю.
4) Верно: дифракция хорошо выражена при размерах препятствия \(d\lesssim\lambda\).
5) Верно: согласно планетарной модели в центре атома расположено положительно заряженное ядро; ответ: \(1,\ 4,\ 5\).
}

% Source 4B89D5: https://ege.fipi.ru/bank/questions.php?search=1&proj=BA1F39653304A5B041B656915DC36B38&qid=4B89D5
\BankTask{4B89D5}{
\begin{enumerate}
\item[1)] Энергия упруго деформированной пружины прямо пропорциональна её удлинению.
\item[2)] Теплопередача путём электромагнитного излучения возможна только в атмосфере Земли и не наблюдается в вакууме.
\item[3)] При соединении двух разноимённо заряженных металлических шаров металлической проволокой перераспределение зарядов будет происходить до полного выравнивания потенциалов шаров.
\item[4)] Электромагнитные волны ультрафиолетового диапазона имеют бóльшую длину волны, чем радиоволны.
\item[5)] В нейтральном атоме суммарное количество электронов равно суммарному количеству протонов в ядре этого атома.
\end{enumerate}
}{
\smallskip\noindent\textbf{Решение от GPT}\par
1) Неверно: энергия пружины \(E=\frac{kx^2}{2}\) пропорциональна квадрату удлинения. 2) Неверно: электромагнитное излучение переносит энергию и в вакууме. 3) Верно: перераспределение зарядов прекращается при равенстве потенциалов \(\varphi_1=\varphi_2\). 4) Неверно: длина волны ультрафиолетового излучения меньше длины радиоволн. 5) Верно: нейтральность атома означает \(N_e=N_p\). Ответ: \(35\).
}

% Source 9E0663: https://ege.fipi.ru/bank/questions.php?search=1&proj=BA1F39653304A5B041B656915DC36B38&qid=9E0663
\BankTask{9E0663}{
\begin{enumerate}
\item[1)] Если в данной системе отсчёта материальная точка совершила ровно один полный оборот по окружности, то перемещение материальной точки в этой системе отсчёта равно нулю.
\item[2)] Пар над поверхностью жидкости является насыщенным, если за одно и то же время с поверхности жидкости в среднем вылетает меньшее число молекул, чем возвращается обратно в жидкость.
\item[3)] Разноимённые точечные электрические заряды притягиваются друг к другу.
\item[4)] Индукционный ток возникает в замкнутом проводящем контуре при изменении магнитного потока, пронизывающего поверхность, ограниченную контуром.
\item[5)] Термоядерными реакциями называют экзотермические реакции распада тяжёлых ядер.
\end{enumerate}
}{
\smallskip\noindent\textbf{Решение от GPT}\par
1) Верно: после полного оборота конечное положение совпадает с начальным, поэтому \(\Delta\vec r=0\). 2) Неверно: насыщенный пар находится в динамическом равновесии с жидкостью, и средние числа вылетевших и вернувшихся молекул за одинаковое время равны. 3) Верно: разноимённые заряды притягиваются по закону Кулона. 4) Верно: по закону электромагнитной индукции \(\mathcal{E}_i=-\frac{d\Phi}{dt}\), и в замкнутом проводящем контуре возникает ток. 5) Неверно: термоядерные реакции представляют собой слияние лёгких ядер при высоких температурах. Ответ: \(134\).
}

% Source E70C7D: https://ege.fipi.ru/bank/questions.php?search=1&proj=BA1F39653304A5B041B656915DC36B38&qid=E70C7D
\BankTask{E70C7D}{
\begin{enumerate}
\item[1)] Сила трения скольжения всегда противоположна направлению движения тела относительно опоры.
\item[2)] Внутренняя энергия постоянной массы идеального газа прямо пропорциональна его абсолютной температуре.
\item[3)] Напряжённость поля неподвижного точечного заряда в данной точке обратно пропорциональна расстоянию от этой точки до заряда.
\item[4)] При интерференции когерентных волн от точечного источника максимум наблюдается в том случае, если разность хода двух волн, возбуждающих колебания в этой точке, равна целому числу длин волн.
\item[5)] При \(\alpha\)-распаде заряд ядра увеличивается на два элементарных заряда.
\end{enumerate}
}{
\smallskip\noindent\textbf{Решение от GPT}\par
1) Верно: сила трения скольжения направлена против скорости движения тела относительно опоры. 2) Верно: для постоянной массы идеального газа \(U=\frac{i}{2}\nu RT\), поэтому \(U\propto T\). 3) Неверно: \(E=\frac{k|q|}{r^2}\), то есть напряжённость обратно пропорциональна квадрату расстояния. 4) Верно: условие интерференционного максимума \(\Delta=m\lambda\), где \(m\in\mathbb{Z}\). 5) Неверно: при \(\alpha\)-распаде ядро испускает частицу с зарядом \(+2e\), поэтому заряд оставшегося ядра уменьшается на \(2e\). Ответ: \(124\).
}

% Source 2449F3: https://ege.fipi.ru/bank/questions.php?search=1&proj=BA1F39653304A5B041B656915DC36B38&qid=2449F3
\BankTask{2449F3}{
1) Сила Архимеда, действующая на однородное тело, полностью погружённое в жидкость, прямо пропорциональна объёму тела.

2) При изотермическом расширении внутренняя энергия постоянной массы одноатомного идеального газа уменьшается.

3) Модуль сил взаимодействия двух неподвижных точечных заряженных тел в вакууме обратно пропорционален произведению модулей зарядов.

4) Скорость распространения ультрафиолетового излучения в вакууме равна скорости света в вакууме.

5) В результате электронного \(\beta\)-распада образовавшийся новый элемент смещается на одну клетку ближе к концу Периодической системы элементов Д.\,И. Менделеева.
}{
\smallskip\noindent\textbf{Решение от GPT}\par
1) Верно: для полностью погружённого тела \(F_{\text{А}}=\rho_{\text{ж}}gV\), поэтому сила прямо пропорциональна объёму тела. 2) Неверно: \(U=\frac32\nu RT\), а при изотермическом процессе температура и внутренняя энергия постоянны. 3) Неверно: по закону Кулона \(F=k\frac{|q_1q_2|}{r^2}\), то есть сила прямо пропорциональна произведению модулей зарядов. 4) Верно: ультрафиолетовое излучение является электромагнитным и распространяется в вакууме со скоростью \(c\). 5) Верно: при электронном \(\beta\)-распаде \(Z\) увеличивается на единицу, поэтому элемент смещается на одну клетку к концу таблицы. Ответ: \(1,\ 4,\ 5\).
}

% Source 8213F0: https://ege.fipi.ru/bank/questions.php?search=1&proj=BA1F39653304A5B041B656915DC36B38&qid=8213F0
\BankTask{8213F0}{
1) При равномерном прямолинейном движении за любые равные промежутки времени тело совершает одинаковые перемещения.

2) Средняя кинетическая энергия поступательного теплового движения молекул газа обратно пропорциональна абсолютной температуре газа.

3) В однородном электростатическом поле работа по перемещению заряда между двумя точками не зависит от траектории.

4) При переходе электромагнитной волны из оптически менее плотной в оптически более плотную среду частота волны уменьшается.

5) При электронном \(\beta\)-распаде масса ядра остаётся практически неизменной.
}{
\smallskip\noindent\textbf{Решение от GPT}\par
1) Верно: \(\Delta\vec r=\vec v\Delta t\), и при постоянной скорости равным промежуткам времени соответствуют одинаковые перемещения. 2) Неверно: \(\langle E_{\text{к}}\rangle=\frac32 k_{\text{B}}T\), то есть энергия прямо пропорциональна абсолютной температуре. 3) Верно: электростатическое поле потенциально, и работа \(A=q(\varphi_1-\varphi_2)\) зависит только от начальной и конечной точек. 4) Неверно: при переходе между средами частота сохраняется, изменяются скорость и длина волны. 5) Верно: при электронном \(\beta\)-распаде число нуклонов сохраняется, поэтому масса ядра меняется лишь незначительно. Ответ: \(1,\ 3,\ 5\).
}

\end{document}
